\documentclass[11pt,letterpaper]{article}
\usepackage[T1]{fontenc}
\usepackage[utf8]{inputenc}
\usepackage{lmodern}
\usepackage[margin=1in,headheight=15pt]{geometry}
\usepackage{amsmath,amssymb,amsthm,mathtools,mathrsfs}
\usepackage{microtype,booktabs,longtable,array,enumitem}
\usepackage[dvipsnames]{xcolor}
\usepackage{fancyhdr,aliascnt}
\usepackage{xurl}
\usepackage[colorlinks=true,linkcolor=MidnightBlue,citecolor=MidnightBlue,urlcolor=MidnightBlue,breaklinks=true]{hyperref}
\usepackage[nameinlink,noabbrev]{cleveref}
\hypersetup{pdftitle={Nonabelian Support Obstructions in Surcomplex Analysis},pdfauthor={Research manuscript},pdfsubject={Hahn-coherent matrix Cousin problems and integral vector bundles}}
\pagestyle{fancy}
\fancyhf{}
\fancyhead[L]{\small Nonabelian support obstructions}
\fancyhead[R]{\small Surcomplex analysis}
\fancyfoot[C]{\thepage}
\setlength{\parskip}{0.28em}
\setlength{\parindent}{1.2em}
\setlist{nosep,leftmargin=2em}
\emergencystretch=2em
\newtheorem{theorem}{Theorem}[section]
\newaliascnt{lemma}{theorem}
\newtheorem{lemma}[lemma]{Lemma}
\aliascntresetthe{lemma}
\newaliascnt{proposition}{theorem}
\newtheorem{proposition}[proposition]{Proposition}
\aliascntresetthe{proposition}
\newaliascnt{corollary}{theorem}
\newtheorem{corollary}[corollary]{Corollary}
\aliascntresetthe{corollary}
\theoremstyle{definition}
\newaliascnt{definition}{theorem}
\newtheorem{definition}[definition]{Definition}
\aliascntresetthe{definition}
\newaliascnt{example}{theorem}
\newtheorem{example}[example]{Example}
\aliascntresetthe{example}
\theoremstyle{remark}
\newaliascnt{remark}{theorem}
\newtheorem{remark}[remark]{Remark}
\aliascntresetthe{remark}
\newcommand{\C}{\mathbb C}
\newcommand{\R}{\mathbb R}
\newcommand{\Q}{\mathbb Q}
\newcommand{\Z}{\mathbb Z}
\newcommand{\N}{\mathbb N}
\newcommand{\No}{\mathbf{No}}
\newcommand{\OO}{\mathcal O}
\newcommand{\HH}{\mathscr H_\Gamma}
\newcommand{\II}{\mathscr I_\Gamma}
\newcommand{\JJ}{\mathscr J_\Gamma}
\newcommand{\GG}{\mathscr G_{r,\Gamma}}
\newcommand{\Mat}{\operatorname{Mat}}
\newcommand{\GL}{\operatorname{GL}}
\newcommand{\SL}{\operatorname{SL}}
\newcommand{\supp}{\operatorname{supp}}
\newcommand{\red}{\operatorname{red}}
\newcommand{\Pol}{\operatorname{Pol}}
\newcommand{\Reg}{\operatorname{Reg}}
\newcommand{\End}{\mathcal E nd}
\newcommand{\id}{\operatorname{id}}
\newcommand{\PP}{\mathcal P}
\newcommand{\EU}{\mathcal E}
\newcommand{\sfam}{\mathcal S}
\newcommand{\angles}[1]{\langle #1\rangle_{\N}}
\newcolumntype{L}[1]{>{\raggedright\arraybackslash}p{#1}}
\numberwithin{equation}{section}

\begin{document}
\hypersetup{pageanchor=false}
\begin{titlepage}
\centering
\vspace*{0.5cm}
{\large\scshape A research continuation of the Surreal repository\par}
\vspace{0.8cm}
{\Huge\bfseries Nonabelian Support\par Obstructions in\par Surcomplex Analysis\par}
\vspace{0.6cm}
{\LARGE A Sharp Matrix Cousin Criterion\par and Arbitrarily Deep Unipotent Examples\par}
\vspace{0.75cm}
{\large September 21, 2026\par}
\vspace{0.65cm}
\begin{minipage}{0.94\textwidth}
\begin{abstract}
We study holomorphic matrix gluing over the integral Hahn coefficient sheaf
$\II=\OO((t^\Gamma))_{\geq0}$, with transition matrices congruent to the identity modulo positive exponents. Here $\Gamma$ is a set-sized ordered abelian group; when $\Gamma\subseteq\No$, the construction has a canonical surcomplex interpretation. On a puncture-and-disk cover of an open Riemann surface, we prove an exact global splitting criterion: independently factor each transition matrix into its normalized polar and holomorphic factors, and test whether the union of the supports of the polar factors is well ordered. This criterion is invariant under arbitrary local holomorphic right gauges and under ordered value-group extension.

The criterion produces genuinely nonabelian counterexamples. In rank three, all original nonremovable matrix coefficients use a single Hahn exponent, the determinant is one, and the abelianized gluing problem splits; nevertheless the integral vector bundle is nontrivial. A missing matrix cross term creates the obstruction, while inserting that term removes it even though it introduces descending exponents among the original poles. In every rank $r\geq3$ we construct an example whose quotient by the last central subgroup of the upper-unitriangular group is trivial, but whose full integral vector bundle is not. Thus support obstructions occur at arbitrarily deep nilpotent levels.

We also prove support-controlled rigidity for positive matrix cocycles on Stein manifolds, and for positive deformations of ordinary holomorphic vector bundles. On Stein bases, every positive cocycle over a cyclic value group splits; noncyclic groups admit determinant-one counterexamples already in rank two. All infinite constructions use well-ordered supports and finite coefficient multiplicity, not fine-topological convergence.
\end{abstract}
\end{minipage}
\vfill
\begin{minipage}{0.94\textwidth}\small
\textbf{Status and scope.} The exact global matrix criterion and the explicit deep support obstructions are proposed new results in this specified Hahn-coherent framework. The local factorization mechanism and the ordinary analytic inputs are established mathematics. A targeted literature check did not locate the global criterion or these examples; this is not a certification of priority. The article gives proofs, not a claim to have resolved a named published conjecture. The proofs are not refereed or machine checked. Nontriviality is proved over the \emph{integral} coefficient sheaf; no unproved passage to the sheaf obtained by inverting all monomials is made.
\end{minipage}
\end{titlepage}
\hypersetup{pageanchor=true}
\pagenumbering{roman}
\tableofcontents
\newpage
\pagenumbering{arabic}

\section{The problem and the contribution}\label{sec:introduction}

\subsection{The gap being addressed}

The repository \texttt{VladimirReshetnikov/Surreal} develops several distinct analytic structures over surreal and surcomplex numbers. Its global-divisor report proves scalar support criteria and a Picard-group dichotomy, then explicitly leaves finite-rank vector bundles as a further direction: matrix logarithms do not turn multiplication into addition, so a noncommutative support construction is needed. We take that direction rather than repeat the existing scalar, local-zero, contour, differential-algebra, or rank-one results. The repository review is pinned to commit
\begin{center}\small\ttfamily
39f2be6667ade51bca2b45daa47e289d69c09764.
\end{center}
The relevant source is the subsection ``Further precise directions'' in the global-divisor article \cite{RepoGlobal}; the catalogue's distinctions between coefficient rings remain essential \cite{RepoMap}.

The basic question is concrete. Let $A\subset\C$ be closed and discrete, let $D_a$ be pairwise disjoint ordinary disks about $a\in A$, and put
\[
 U_0=\C\setminus A,\qquad W_a=D_a\setminus\{a\}.
\]
Suppose that, separately for each $a$, we have a legitimate positive Hahn matrix
\begin{equation}\label{eq:inputintro}
 G_a=I+\sum_{\gamma>0}G_{a,\gamma}(z)t^\gamma,
 \qquad G_{a,\gamma}\in\Mat_r(\OO(W_a)),
\end{equation}
with well-ordered support. There need not be one well-ordered set containing the supports for all $a$. When can one find positive matrices $B_0$ on $U_0$ and $B_a$ on $D_a$ such that
\begin{equation}\label{eq:splitintro}
 G_a=B_0B_a^{-1}\quad\text{on }W_a?
\end{equation}
Each local equation is harmless. The difficulty is the requirement that the entries of the \emph{one global matrix} $B_0$ have legitimate Hahn supports.

For scalar matrices, a logarithm reduces the problem to additive principal parts. For matrices of size at least two,
\[
 \log(AB)\ne\log A+\log B
\]
in general. Discarding the cross terms can change the answer to the existence problem. Our rank-three example makes that failure explicit, without a norm estimate, a small-divisor issue, or a delicate analytic singularity.

\subsection{The principal results}

Every matrix \eqref{eq:inputintro} admits a unique local factorization
\begin{equation}\label{eq:factorintro}
 G_a=P_aH_a,
\end{equation}
where $H_a$ is positive and holomorphic on $D_a$, while every nonconstant Hahn coefficient of $P_a$ is a negative Laurent part in $u=z-a$. Equivalently, that coefficient extends to $\C\setminus\{0\}$ and vanishes at infinity in the auxiliary $u$-plane. This is a normalized Birkhoff-type factorization, not the Hilbert-space polar decomposition of a matrix.

The first main result is the equivalence
\begin{equation}\label{eq:criterionintro}
 \boxed{\quad
 \eqref{eq:splitintro}\text{ is solvable}
 \quad\Longleftrightarrow\quad
 \bigcup_{a\in A}\supp(P_a-I)\text{ is well ordered}.
 \quad}
\end{equation}
The support is the union across all matrix entries. If that set is $E$, the central matrix $B_0-I$ can be supported inside $\angles E\setminus\{0\}$, the positive part of the additive monoid generated by $E$. The disk matrices may need the original, individually admissible supports as well. The distinction is necessary: freely removable local terms should not be charged to the global obstruction.

The second result is a pair of opposite examples. Over $\Gamma=\Q$, let $a=n\geq2$, $u=z-n$, and write $E_{ij}$ for matrix units. The matrices
\begin{equation}\label{eq:rank3intro}
 G_n=I+\frac{t}{u}E_{12}+t^{1/n}E_{23}
\end{equation}
have determinant one and no original singular coefficient except the exponent-$1$ coefficient. Nevertheless their normalized polar factors are
\begin{equation}\label{eq:rank3polarintro}
 P_n=I+\frac{t}{u}E_{12}-\frac{t^{1+1/n}}{u}E_{13}.
\end{equation}
The descending set $\{1+1/n:n\geq2\}$ prevents a splitting. Adding $t^{1+1/n}E_{13}/u$ to \eqref{eq:rank3intro} makes the problem split. Thus the support of the raw singular coefficients is neither a sufficient nor a necessary test.

The third result places the obstruction at any prescribed finite nilpotent depth. For every $r\geq3$ there is an upper-unitriangular positive cocycle on the plane which becomes trivial modulo the last central subgroup
\[
 Z_r=\{I+cE_{1r}\},
\]
but remains nontrivial even when a general positive $r\times r$ matrix trivialization is allowed. Its integral vector bundle is nontrivial, has trivial reduction and trivial determinant, and is trivial on every bounded ordinary domain. These conclusions are not merely failures to preserve a selected upper-triangular flag.

\subsection{What is and is not being claimed as new}

Hahn support calculus is classical \cite{Neumann}. The extension of restricted analytic functions to surreal workspaces is classical \cite{vdDE}. Algebraic Birkhoff factorization in complete filtered noncommutative algebras, including its recursive form and the necessity of commutator corrections, is also established \cite{EFGK,EFMP}. We reprove the needed local factorization with explicit arbitrary-order-group support control instead of claiming a new factorization principle.

The proposed contribution is the \emph{exact global support invariant} for an infinite family of matrix gluing data, its relationship to actual integral vector-bundle triviality, and the explicit hierarchy of nonabelian obstructions. The Stein rigidity statement supplies useful context and a positive counterpart; its mechanism is the standard deformation-theoretic use of ordinary coherent-sheaf acyclicity, with the Hahn support issue made explicit. The rank-one specialization recovers, rather than replaces or newly proves priority for, the repository's scalar criterion.

We do not classify all bundles over every ordinary cover. We solve the puncture-and-disk problem for arbitrary finite matrix rank, and prove a sufficient theorem for arbitrary Stein covers with globally admissible transition support. We also carefully distinguish the integral coefficient sheaf from its monomial localization. These boundaries are part of the theorem statements, not qualifications added after their proofs.

\section{The coefficient category and its support calculus}\label{sec:coefficients}

\subsection{One set-sized workspace}

Let $\Gamma$ be a set-sized totally ordered abelian group, written additively. For a complex algebra $R$, the Hahn algebra $R((t^\Gamma))$ consists of formal sums
\[
 F=\sum_{\gamma\in S}f_\gamma t^\gamma,
 \qquad f_\gamma\in R,
\]
whose nonzero support $S$ is well ordered in the given order. Addition is coefficientwise and multiplication is convolution. For matrices we use the union of the supports of their finitely many entries. An exponent occurs in that union exactly when the corresponding coefficient matrix is not identically zero.

When $\Gamma\subseteq\No$ is an ordered additive subgroup, Conway normal forms give the interpretation
\begin{equation}\label{eq:surrealembedding}
 \C((t^\Gamma))\longrightarrow\No[i],
 \qquad \sum c_\gamma t^\gamma\longmapsto
 \sum c_\gamma\omega^{-\gamma}.
\end{equation}
Real and imaginary parts are interpreted separately. All supports here are sets. No sum indexed by a proper class is needed. Divisibility of $\Gamma$ is unnecessary for our arguments because we do not use algebraic closure or extraction of roots.

All ordinary manifolds and Riemann surfaces have their usual Hausdorff, second-countable topology. For a connected ordinary complex domain $V$, set
\begin{align}
 \HH(V)&=\OO(V)((t^\Gamma)),\label{eq:H}\\
 \II(V)&=\{F\in\HH(V):\supp F\subseteq\Gamma_{\geq0}\},\label{eq:I}\\
 \JJ(V)&=\{F\in\HH(V):\supp F\subseteq\Gamma_{>0}\}.\label{eq:J}
\end{align}
Every coefficient must be holomorphic on the \emph{same fixed domain} $V$. We do not substitute a family of convergent germs with unrelated radii. For disconnected open sets, sections are specified component by component. This is the sheaf convention throughout.

There is a reduction map
\begin{equation}\label{eq:reduction}
 \red:\II\longrightarrow\OO,\qquad
 \sum_{\gamma\geq0}f_\gamma t^\gamma\longmapsto f_0,
 \qquad\ker\red=\JJ.
\end{equation}
Define the positive matrix group
\begin{equation}\label{eq:G}
 \GG(V)=I+\Mat_r(\JJ(V)).
\end{equation}
Its elements are automatically invertible, as proved below. ``Positive'' refers to Hahn exponents, not positivity of a Hermitian matrix or of its entries.

\subsection{Why this is a sheaf on the ordinary base}

There is a minor but important support point in gluing. A nonzero ordinary holomorphic coefficient on a connected domain cannot vanish on a nonempty open subset. Thus restriction of a Hahn series between nonempty connected ordinary domains preserves its support. If compatible local sections are given on a cover of a connected domain, the identity theorem along chains of intersecting patches shows that their coefficient functions glue and that the global support equals the support seen on any one nonempty connected patch. This support is well ordered. For disconnected opens we take the product over connected components.

In particular, if a section on $X\setminus A$ extends separately across each point of a discrete $A$, it extends with its \emph{existing global support}; no additional exponents can appear only at an isolated ordinary point. This observation is used repeatedly when matrix quotients are extended.

\subsection{The support lemma}

\begin{lemma}[Neumann support calculus]\label{lem:Neumann}
Let $E\subset\Gamma_{>0}$ be well ordered, and let $M=\angles E$ include the empty sum $0$. Then $M$ is well ordered. For each $\gamma\in\Gamma$, only finitely many finite words in $E$ have sum $\gamma$. Consequently, if $S\subset\Gamma$ is well ordered, then $S+M$ is well ordered, and every fixed coefficient in a sum formed by adjoining finitely many positive-support factors receives only finitely many contributions, provided each input family is Hahn summable.
\end{lemma}

The first two assertions are the positive-support form of Neumann's lemma \cite{Neumann}. The remaining assertion combines them with the standard finite-decomposition property for sums of two well-ordered subsets of an ordered abelian group. We use this established combinatorial theorem as an input. It is stronger than saying merely that all exponents are positive.

\begin{remark}[No hidden Archimedean assumption]\label{rem:nonarch}
A least positive input exponent $\delta$ need not have the property that $n\delta$ eventually exceeds a prescribed $\gamma$. For example, in $\Z^2$ with lexicographic order, $n(0,1)<(1,0)$ for every ordinary integer $n$. Our proofs never infer coefficient finiteness from that false cofinality assertion. They use the finite-word conclusion of \cref{lem:Neumann} and recursion on the well-ordered set $M$.
\end{remark}

\begin{lemma}[Positive inversion and formal functions]\label{lem:inverse}
For $A\in\Mat_r(\JJ(V))$, the expressions
\[
 (I+A)^{-1}=\sum_{k\geq0}(-A)^k,
 \qquad
 \exp A=\sum_{k\geq0}\frac{A^k}{k!},
 \qquad
 \log(I+A)=\sum_{k\geq1}\frac{(-1)^{k+1}}{k}A^k
\]
are Hahn summable coefficientwise on $V$. Their nonconstant supports lie in $\angles{\supp A}\setminus\{0\}$. The inverse is two-sided, and exponential and logarithm are mutually inverse as one-variable formal substitutions. They do not make matrix multiplication additive.
\end{lemma}

\begin{proof}
Expand an entry of $A^k$. Its terms are indexed by a word of $k$ positive exponents and a path through $r$ matrix indices. At a fixed total exponent only finitely many exponent words occur by \cref{lem:Neumann}; each word has finitely many matrix-index paths. Hence all displayed coefficient sums are finite and their coefficients are holomorphic on the original $V$. Ordinary formal identities involving one variable, or the geometric-series cancellation, may therefore be checked coefficientwise. Noncommuting independent variables do not satisfy the scalar exponential addition law.
\end{proof}

\subsection{Connection with surcomplex evaluation}

At an ordinary point $c\in V$ and a positive-valuation infinitesimal $\epsilon$, a coefficient function is evaluated by its ordinary Taylor series:
\[
 f_\gamma(c+\epsilon)
 =\sum_{k\geq0}\frac{f_\gamma^{(k)}(c)}{k!}\epsilon^k.
\]
If $S$ is the original Hahn support and $T=\supp\epsilon$, the resulting support is contained in $S+\angles T$. The same support lemma makes the evaluation legitimate, entry by entry. Thus our holomorphic matrices give matrices on finite surcomplex thickenings of ordinary charts. The ordinary complex base is not being replaced by a fine-topological surcomplex manifold. No assertion that ordinary paths are continuous in the fine surreal topology is required.

\section{Ordinary Laurent splitting and analytic interpolation}\label{sec:ordinary}

\subsection{A fixed-disk decomposition}

Fix $D=\{|u|<\rho\}$ and $D^*=D\setminus\{0\}$. A function $f\in\OO(D^*)$ has the Laurent decomposition
\begin{equation}\label{eq:Laurent}
 f=f_-+f_+,
 \qquad
 f_- =\sum_{k\geq1}a_{-k}u^{-k},
 \qquad
 f_+ =\sum_{k\geq0}a_ku^k.
\end{equation}
The negative part converges for every $u\ne0$ and vanishes at $u=\infty$; the positive part is holomorphic on the original disk $D$. Write $\PP$ for the algebra of all such negative parts and write $\pi_-,\pi_+$ for the projections. These act entrywise on matrices.

\begin{lemma}[The ordinary splitting algebras]\label{lem:Laurent}
There is a direct-sum decomposition of complex vector spaces
\[
 \OO(D^*)=\PP\oplus\OO(D).
\]
Both summands are closed under multiplication, although $\PP$ has no multiplicative identity. Their intersection is zero. The decomposition does not shrink $D$.
\end{lemma}

\begin{proof}
For every $s\in(0,\rho)$, Cauchy's formula bounds the negative Laurent coefficients by $|a_{-k}|\leq M_s s^k$. Consequently the series in $1/u$ converges for $|u|>s$, and $s$ can be arbitrarily small. It vanishes at infinity. Conversely, a function holomorphic on $\C\setminus\{0\}$ and vanishing at infinity has only negative Laurent powers. Products have the same properties. A function in the intersection extends holomorphically across zero and across infinity, so it is constant on the Riemann sphere; its value at infinity is zero. The positive Laurent part converges on the original disk of radius $\rho$.
\end{proof}

The negative part may have an essential singularity at zero. We deliberately do not impose a finite pole-order bound. The explicit obstruction examples below use only simple poles.

\subsection{The scalar analytic input}

\begin{lemma}[Isolated singular parts on the plane]\label{lem:MLplane}
Let $A\subset\C$ be closed and discrete. At every $a\in A$, let $p_a\in\PP_a$ be a prescribed negative Laurent part in $z-a$. There exists $f\in\OO(\C\setminus A)$ such that $f-p_a$ extends holomorphically through $D_a$. One may additionally prescribe $f(z_*)$ at one fixed $z_*\notin A$.
\end{lemma}

\begin{proof}
The finite case is a finite sum. In the infinite case enumerate $A=\{a_n\}$ with $|a_n|\to\infty$, and choose radii $R_n\to\infty$ satisfying $R_n<|a_n|$ for all sufficiently large $n$. Each $p_{a_n}(z-a_n)$ is holomorphic on $|z|<|a_n|$. Choose a Taylor polynomial $q_n$ at zero such that
\[
 \sup_{|z|\leq R_n}|p_{a_n}(z-a_n)-q_n(z)|<2^{-n}.
\]
After separating finitely many early indices, the series
\[
 f(z)=\sum_n\bigl(p_{a_n}(z-a_n)-q_n(z)\bigr)
\]
converges normally on every compact subset of $\C\setminus A$. Near a fixed $a_m$, all terms except the $m$th form a normally convergent holomorphic series; hence the required negative Laurent part is $p_{a_m}$. The differences extend over the full $D_{a_m}$ by uniqueness. Adding a constant fixes $f(z_*)$.
\end{proof}

This proof allows essential isolated singularities because it approximates the entire negative part away from its centre. Nothing in it uses a bound on pole orders.

\begin{lemma}[Open-surface version]\label{lem:MLsurface}
The conclusion of \cref{lem:MLplane} holds on any connected open Riemann surface $X$, with a closed discrete $A$, pairwise disjoint coordinate disks $D_a$, and negative Laurent parts in their chosen coordinates.
\end{lemma}

\begin{proof}
An open Riemann surface is Stein by the Behnke--Stein theorem \cite{BS}. The data $p_a$ define an ordinary additive holomorphic cocycle on the cover $X\setminus A,D_a$. Cartan's theorem B gives $H^1(X,\OO)=0$ and therefore a splitting of this cocycle. In the chosen sign convention the central component is the desired $f$. An additive torsor that is trivial has a section expressed in the original charts, so no acyclicity assumption on this particular cover is needed. A constant again prescribes one value. See \cite{Cartan,Schmitt} for the ordinary coherent-sheaf input.
\end{proof}

For the main theorem on the plane, \cref{lem:MLplane} makes the analytic existence step self-contained. The general open-surface statement uses exactly the two established ordinary theorems identified in the proof, not an unproved theorem B for a Hahn sheaf.

\section{Support-controlled local matrix factorization}\label{sec:factorization}

\subsection{The normalized factors}

For the chosen disk and coordinate, let
\[
 \mathscr P_{r,\Gamma}
 =I+\Mat_r\bigl(\PP((t^\Gamma))_{>0}\bigr).
\]
The notation means that the nonconstant coefficients lie in $\PP$ and have positive well-ordered support. This is a group: nonconstant coefficients of a product or a positive inverse are finite sums of products of negative Laurent parts. Each such product still vanishes at auxiliary infinity.

\begin{theorem}[Local normalized factorization]\label{thm:local}
Let $G\in\GG(D^*)$ and let $S=\supp(G-I)$. There exist unique matrices
\begin{equation}\label{eq:localfactor}
 P\in\mathscr P_{r,\Gamma},\qquad H\in\GG(D),
 \qquad G=PH.
\end{equation}
Both nonconstant supports lie in $M=\angles S\setminus\{0\}$. We denote the factors by $P=\Pol(G)$ and $H=\Reg(G)$. Every coefficient of $H$ is holomorphic on the original disk, not on an exponent-dependent smaller disk.
\end{theorem}

\begin{proof}
The case $S=\varnothing$ is immediate. Otherwise write
\[
 G=I+\sum_{\gamma>0}g_\gamma t^\gamma,
 \quad P=I+\sum_{\gamma\in M}p_\gamma t^\gamma,
 \quad H=I+\sum_{\gamma\in M}h_\gamma t^\gamma.
\]
We use well-founded recursion on the positive part $M$ of the generated monoid. Put $g_\gamma=0$ for exponents not originally present. Once all smaller coefficients have been chosen, define
\begin{align}
 r_\gamma
   &=g_\gamma-
       \sum_{\substack{\eta+\theta=\gamma\\\eta,\theta>0}}
         p_\eta h_\theta,\label{eq:factorrecursion1}\\
 p_\gamma&=\pi_-(r_\gamma),\qquad
 h_\gamma=\pi_+(r_\gamma).\label{eq:factorrecursion2}
\end{align}
Every summand involves strictly smaller exponents, and the sum is finite by \cref{lem:Neumann}. Each step therefore produces a negative Laurent matrix and a matrix holomorphic on $D$. The series have supports in $M$ and satisfy $PH=G$ coefficientwise. Their inverses exist by \cref{lem:inverse}.

For uniqueness, suppose $P_1H_1=P_2H_2$. Then
\[
 P_2^{-1}P_1=H_2H_1^{-1}.
\]
The nonconstant coefficients on the left belong to $\PP$, while those on the right are holomorphic on $D$. After placing the finitely many input supports in a common positive generated monoid, coefficientwise comparison and \cref{lem:Laurent} show that both sides equal $I$.
\end{proof}

\begin{proposition}[Right-gauge invariance]\label{prop:rightgauge}
If $Q\in\GG(D)$, then
\begin{equation}\label{eq:rightgauge}
 \Pol(GQ)=\Pol(G).
\end{equation}
The analogous statement is generally false for multiplication on the left by $Q$.
\end{proposition}

\begin{proof}
The factorization $GQ=\Pol(G)(\Reg(G)Q)$ has the defining properties of \cref{thm:local}. Uniqueness proves the equality. Left multiplication can mix regular and polar terms in the opposite order; the examples in \cref{sec:rankthree} show that order matters.
\end{proof}

\subsection{The first terms and the source of the obstruction}

For an ordinary single positive parameter, write
\[
 G=I+t g_1+t^2g_2+t^3g_3+\cdots.
\]
The recursion begins
\begin{align*}
 p_1&=\pi_-g_1,& h_1&=\pi_+g_1,\\
 p_2&=\pi_-(g_2-p_1h_1),&
 h_2&=\pi_+(g_2-p_1h_1),\\
 p_3&=\pi_-(g_3-p_1h_2-p_2h_1).&&
\end{align*}
The second-order correction already contains the product of a polar part and a regular part. It can be nonzero even when $g_2=0$. With different positive exponents at different punctures, the exponents of these products can form a descending sequence that was absent from the raw singular data.

For scalar matrices one has the simpler formula
\begin{equation}\label{eq:scalarPol}
 \Pol(G)=\exp\bigl(\pi_-\log G\bigr),\qquad
 \Reg(G)=\exp\bigl(\pi_+\log G\bigr).
\end{equation}
Here the projections act coefficientwise and the scalar exponential addition law verifies the factorization. In a noncommutative matrix algebra, replacing \cref{thm:local} by \eqref{eq:scalarPol} is invalid. The algebraic Birkhoff literature describes the associated Baker--Campbell--Hausdorff corrections in a filtered setting \cite{EFGK,EFMP}; our direct recursion avoids having to expand that formula.

\subsection{Preservation of triangular shape}

\begin{corollary}[Unitriangular factors]\label{cor:unitriangular}
If $G$ is upper unitriangular, then $\Pol(G)$ and $\Reg(G)$ are upper unitriangular. If a fixed pattern of strictly upper-triangular entries is closed under matrix multiplication, the factors preserve that pattern as well.
\end{corollary}

\begin{proof}
Both Laurent projections preserve zero entries. Every nonlinear term in \eqref{eq:factorrecursion1} is a product of previously constructed strictly upper-triangular matrices and thus stays in the specified multiplicatively closed pattern. Induction proves the claim.
\end{proof}

\section{The sharp global matrix Cousin criterion}\label{sec:global}

Let $X$ be a connected open Riemann surface, let $A\subset X$ be closed and discrete, and choose pairwise disjoint coordinate disks $(D_a,u_a)$ with $u_a(a)=0$. Write
\[
 U_0=X\setminus A,\qquad W_a=D_a\setminus\{a\}.
\]
The punctured surface $U_0$ is connected. For each $a$, independently, let $G_a\in\GG(W_a)$. The $D_a$ do not meet one another, so this is a complete positive matrix cocycle on the displayed cover, with the inverse transition on the reversed overlap.

\begin{definition}[Normalized polar support]\label{def:polsupport}
Factor $G_a=P_aH_a$ by \cref{thm:local}. Define
\begin{equation}\label{eq:polsupport}
 \EU(G)=\bigcup_{a\in A}\supp(P_a-I)\subseteq\Gamma_{>0}.
\end{equation}
The datum is \emph{polar-admissible} when $\EU(G)$ is well ordered. No sum over $a$ is intended in \eqref{eq:polsupport}; it is a set-theoretic union of supports.
\end{definition}

\begin{theorem}[Exact nonabelian Cousin criterion]\label{thm:criterion}
The following statements are equivalent:
\begin{enumerate}[label=\textup{(\roman*)}]
\item There exist $B_0\in\GG(U_0)$ and $B_a\in\GG(D_a)$ satisfying $G_a=B_0B_a^{-1}$ on every $W_a$.
\item The set $\EU(G)$ is well ordered.
\end{enumerate}
When these hold, let $M=\angles{\EU(G)}$. One can choose
\begin{equation}\label{eq:centralsupport}
 \supp(B_0-I)\subseteq M\setminus\{0\}.
\end{equation}
In the local factorizations $B_0=P_aQ_a$, the matrices $Q_a-I$ can also be supported in $M\setminus\{0\}$. The disk trivializers are $B_a=H_a^{-1}Q_a$ and hence have individually legitimate supports. A value $B_0(z_*)=I$ may be imposed at any chosen $z_*\in U_0$.
\end{theorem}

\begin{proof}
\emph{Necessity.} Suppose $G_a=B_0B_a^{-1}$. By right-gauge invariance,
\[
 P_a=\Pol(G_a)=\Pol(B_0|_{W_a}).
\]
Let $S=\supp(B_0-I)$. The local support estimate in \cref{thm:local} is independent of the puncture and of the Laurent coefficients of the input. It gives
\[
 \supp(P_a-I)\subseteq\angles S\setminus\{0\}
 \quad\text{for every }a.
\]
This one generated monoid is well ordered. Its subset $\EU(G)$ is therefore well ordered.

\emph{Sufficiency.} Assume $E=\EU(G)$ is well ordered and set $M=\angles E$. We construct a global $F\in\GG(U_0)$ and local $Q_a\in\GG(D_a)$ satisfying
\begin{equation}\label{eq:Fgoal}
 F=P_aQ_a
\end{equation}
with every nonconstant coefficient supported in $M$. Write
\[
 P_a=I+\sum_{\gamma>0}p_{a,\gamma}t^\gamma,
 \quad F=I+\sum_{\gamma>0}f_\gamma t^\gamma,
 \quad Q_a=I+\sum_{\gamma>0}q_{a,\gamma}t^\gamma.
\]
At a positive exponent $\gamma\in M$, suppose all smaller coefficients have been constructed. Define on $W_a$
\begin{equation}\label{eq:globalrecursion}
 R_{a,\gamma}
 =p_{a,\gamma}+
   \sum_{\substack{\eta+\theta=\gamma\\\eta,\theta>0}}
      p_{a,\eta}q_{a,\theta}.
\end{equation}
This is a finite sum. Apply \cref{lem:MLsurface}, entry by entry, to the prescribed ordinary negative Laurent matrices $\pi_-R_{a,\gamma}$. It supplies a single matrix
\[
 f_\gamma\in\Mat_r(\OO(U_0))
\]
such that $f_\gamma-R_{a,\gamma}$ extends holomorphically across $D_a$ for every $a$. Set
\begin{equation}\label{eq:qrecursion}
 q_{a,\gamma}=f_\gamma-R_{a,\gamma}.
\end{equation}
The prescribed zero value at $z_*$ can be imposed on every $f_\gamma$.

Recursion on the well-ordered set $M\setminus\{0\}$ constructs all coefficients. The scalar analytic step changes coefficient functions but introduces no new Hahn exponents. Every coefficient of every $Q_a$ is holomorphic on its original disk. The resulting Hahn matrices satisfy \eqref{eq:Fgoal} coefficientwise and are invertible by \cref{lem:inverse}.

Finally $G_a=P_aH_a$ and $F=P_aQ_a$ imply
\[
 G_a=F\,(H_a^{-1}Q_a)^{-1}.
\]
Take $B_0=F$ and $B_a=H_a^{-1}Q_a$. The support of $H_a-I$ lies in the monoid generated by the original support of $G_a-I$; the support of $Q_a-I$ lies in $M$. A finite union of these two well-ordered positive sets generates a legitimate monoid, so $B_a$ is a well-defined section on $D_a$. There is no requirement that the disk supports have a globally well-ordered union.
\end{proof}

\subsection{Why the converse is sharp}

An infinite set of punctures does not itself prevent a solution. At each fixed exponent, the ordinary analytic problem is always solvable. Nor is the obstruction a divergent ordinary sum: the ordinary coefficient construction uses normal convergence, which is independent of the formal Hahn exponents. The failure occurs precisely when the \emph{necessary normalized polar information} cannot fit into a single well-ordered support.

The test cannot be replaced by $\bigcup_a\supp(G_a-I)$. Indeed, $G_a$ may itself be holomorphic on $D_a$, in which case $P_a=I$ and the solution $B_0=I$, $B_a=G_a^{-1}$ exists even if those raw supports have a descending union. More subtly, even discarding all raw holomorphic coefficients is not enough: noncommutative products between discarded and retained terms can create essential polar information. That is the content of \cref{sec:rankthree}.

\subsection{The space of solutions}

\begin{proposition}[Right torsor of splittings]\label{prop:torsor}
For a polar-admissible datum, all splittings are obtained from any one splitting by
\[
 (B_0,B_a)\longmapsto(B_0T,B_aT|_{D_a}),
 \qquad T\in\GG(X).
\]
The action is free and transitive. Under the normalization $B_0(z_*)=I$, use the subgroup $T(z_*)=I$.
\end{proposition}

\begin{proof}
Right multiplication by a common $T$ preserves $B_0B_a^{-1}$. Conversely, from two splittings obtain
\[
 B_0^{-1}B'_0=B_a^{-1}B'_a\quad\text{on }W_a.
\]
The central quotient extends through every disk by the displayed equality. The sheaf observation in \cref{sec:coefficients} gives a global positive section $T$ with one legitimate support on connected $X$. This is the unique element relating the splittings.
\end{proof}

\section{Invariance and the scalar comparison}\label{sec:invariance}

\subsection{Independent local gauges}

\begin{corollary}[Removal of arbitrarily supported regular data]\label{cor:localgauges}
For independently chosen $R_a\in\GG(D_a)$, the families $(G_a)$ and $(G_aR_a)$ have the same normalized polar factors and the same solvability status. No well-ordering condition on $\bigcup_a\supp(R_a-I)$ is needed.
\end{corollary}

This is stronger than invariance under finitely many transformations. The infinitely many regular parts are genuinely irrelevant \emph{after} they have been removed in the correct multiplication order.

\subsection{Changing coordinates or workspace}

\begin{corollary}[Coordinate independence of admissibility]\label{cor:coordinates}
Changing the local coordinates, or shrinking the disks, can change individual normalized polar coefficients but cannot change whether their global support union is well ordered.
\end{corollary}

\begin{proof}
Existence of the splitting in \cref{thm:criterion} is invariant under these changes. Shrinking a disk cannot turn a previously nonextendible germ into an extendible one; a local splitting on smaller disks yields extension of the same central matrix and hence a splitting on the original disks by analytic continuation of $G_a^{-1}B_0$. Alternatively, reexpressing a previously constructed polar factor in the new coordinate and applying \cref{thm:local} puts every new support in the monoid generated by the old support union. The inverse coordinate change gives the converse.
\end{proof}

\begin{corollary}[Absoluteness under ordered extension]\label{cor:extension}
Let $\Gamma\hookrightarrow\Delta$ be an order-preserving embedding of set-sized ordered abelian groups. A datum defined over $\Gamma$ splits in the positive matrix category over $\Delta$ if and only if it splits over $\Gamma$.
\end{corollary}

\begin{proof}
The recursive local polar factor computed in $\Gamma$ is also a factor in $\Delta$ and remains the unique normalized factor. Thus its support union is the image of $\EU(G)$. An order embedding preserves and reflects well ordering of a subset. Apply \cref{thm:criterion} in both groups.
\end{proof}

In particular, the obstructions below cannot be removed merely by allowing more surreal exponents in the same integral positive-matrix problem. The assertion is not a claim about a change to another analytic category or about allowing negative-valuation gauges.

\subsection{Recovery of the rank-one singular-support test}

For $r=1$, let $L_a=\log G_a$ and define the logarithmic singular support
\[
 \Sigma(G)=\bigcup_a\supp(\pi_-L_a).
\]
The formula \eqref{eq:scalarPol} gives
\begin{equation}\label{eq:scalarequiv}
 \EU(G)\text{ well ordered}
 \quad\Longleftrightarrow\quad
 \Sigma(G)\text{ well ordered}.
\end{equation}
For the forward direction, $\log P_a=\pi_-L_a$ is supported in the positive monoid generated by $\EU(G)$. For the reverse direction, exponentiation puts every $P_a-I$ inside the monoid generated by $\Sigma(G)$. This is exactly the scalar logarithmic mechanism already used in \cite{RepoGlobal}. No analogous entrywise formula for matrix logarithms is substituted for the normalized polar invariant.

\section{A rank-three obstruction invisible in the raw poles}\label{sec:rankthree}

\subsection{The closed factorization in the Heisenberg group}

Let
\[
 G=\begin{pmatrix}1&a&b\\0&1&c\\0&0&1\end{pmatrix},
 \qquad a,b,c\in\JJ(D^*).
\]
Write $a=a_-+a_+$ and $c=c_-+c_+$ by coefficientwise Laurent projection. Direct multiplication gives
\begin{equation}\label{eq:HeisenbergP}
 \Pol(G)=
 \begin{pmatrix}
  1&a_-&\pi_-(b-a_-c_+)\\
  0&1&c_-\\
  0&0&1
 \end{pmatrix},
\end{equation}
with holomorphic factor
\begin{equation}\label{eq:HeisenbergH}
 \Reg(G)=
 \begin{pmatrix}
  1&a_+&\pi_+(b-a_-c_+)\\
  0&1&c_+\\
  0&0&1
 \end{pmatrix}.
\end{equation}
All products are legitimate because this is a finite expression in finitely many Hahn series. Uniqueness in \cref{thm:local} proves that these are the canonical factors. The subtraction of $a_-c_+$ is the essential term.

\subsection{The obstructed datum}

Take $X=\C$, $A=\{n:n\geq2\}$, $D_n=\{|z-n|<1/4\}$, and $\Gamma=\Q$. Put $u_n=z-n$. Consider
\begin{equation}\label{eq:badG}
 G_n=I+\frac{t}{u_n}E_{12}+t^{1/n}E_{23}.
\end{equation}
Each matrix is a finite Hahn expression, has determinant one, and reduces to $I$. Equation \eqref{eq:HeisenbergP} yields
\begin{equation}\label{eq:badPH}
 P_n=I+\frac{t}{u_n}E_{12}-\frac{t^{1+1/n}}{u_n}E_{13},
 \qquad H_n=I+t^{1/n}E_{23}.
\end{equation}
The product $P_nH_n$ has zero $(1,3)$ entry because its two terms cancel.

\begin{theorem}[Hidden rank-three obstruction]\label{thm:hidden}
The matrices \eqref{eq:badG} do not admit a splitting by general positive $3\times3$ matrices. Their abelianized upper-unitriangular cocycle does split. Their raw nonremovable coefficient support is the singleton $\{1\}$, and their determinant cocycle is identically one.
\end{theorem}

\begin{proof}
The normalized polar support contains
\[
 \{1+1/n:n\geq2\},
\]
a strictly decreasing infinite sequence. Hence it is not well ordered, and \cref{thm:criterion} rules out a splitting even in the full positive matrix group.

The abelianization of the upper-unitriangular $3\times3$ group remembers the $(1,2)$ and $(2,3)$ entries and adds them componentwise. Let $\phi\in\OO(U_0)$ have principal part $1/(z-n)$ at every $n\geq2$, with no singularities outside $A$. Such a $\phi$ exists by \cref{lem:MLplane}. One explicit choice is
\begin{equation}\label{eq:phi}
 \phi(z)=\sum_{n=2}^\infty\left(\frac1{z-n}+\frac1n\right),
\end{equation}
whose $n$th summand is $O_K(n^{-2})$ on each compact set away from $A$. The $(1,2)$ component splits using the central function $t\phi$ and the disk functions $t(\phi-1/u_n)$. The $(2,3)$ component splits using central function $0$ and disk functions $-t^{1/n}$. Finally, the only original singular coefficient is $E_{12}/u_n$ at exponent $1$, and unitriangular matrices have determinant one.
\end{proof}

The example is stronger than a determinant obstruction. It is stronger even than the statement that a noncommutative cocycle has nontrivial abelianization. The abelianized problem is solvable; the obstruction appears when its solutions are required to fit together in the original multiplication law.

\subsection{Adding a singular term repairs the problem}

Now define
\begin{equation}\label{eq:goodG}
 \widetilde G_n
 =I+\frac{t}{u_n}E_{12}+t^{1/n}E_{23}
      +\frac{t^{1+1/n}}{u_n}E_{13}.
\end{equation}
Then
\[
 \widetilde G_n
 =\left(I+\frac{t}{u_n}E_{12}\right)
  \left(I+t^{1/n}E_{23}\right),
\]
so $\Pol(\widetilde G_n)=I+tE_{12}/u_n$ and $\EU(\widetilde G)=\{1\}$. A global splitting exists.

For an explicit central matrix take $F=I+t\phi E_{12}$ with \eqref{eq:phi}. Put
\[
 Q_n=I+t\left(\phi-\frac1{u_n}\right)E_{12},
 \qquad \widetilde H_n=I+t^{1/n}E_{23}.
\]
Then $F=\Pol(\widetilde G_n)Q_n$ and
\[
 \widetilde G_n=F\,(\widetilde H_n^{-1}Q_n)^{-1}.
\]
All disk factors are holomorphic on $D_n$. Yet the raw singular support of \eqref{eq:goodG} contains the descending set $\{1+1/n\}$. This proves the other half of the claim that raw singular support is the wrong invariant.

\subsection{Why an entrywise logarithm is not a repair}

The actual matrix logarithm of \eqref{eq:badG} is finite:
\[
 \log G_n=
 \frac{t}{u_n}E_{12}+t^{1/n}E_{23}
 -\frac12\frac{t^{1+1/n}}{u_n}E_{13}.
\]
Exponentiating the negative Laurent part alone would give a $(1,3)$ coefficient $-\tfrac12t^{1+1/n}/u_n$, not the correct coefficient $-t^{1+1/n}/u_n$ in \eqref{eq:badPH}. Thus even after using the true matrix logarithm, one still needs the noncommutative factorization correction. The global criterion is built from the actual normalized factor, not from an ad hoc linearization.

\section{Obstructions at arbitrary nilpotent depth}\label{sec:depth}

\subsection{A chain identity}

Let $r\geq3$, let $a$ be a negative Laurent Hahn function with positive support, and let $c_2,\ldots,c_{r-1}$ be positive Hahn constants. Define
\[
 C=\sum_{j=2}^{r-1}c_jE_{j,j+1},\qquad
 H=I+C,\qquad G=I+C+aE_{12}.
\]
Since $C^{r-1}=0$,
\begin{equation}\label{eq:chainP}
 P=GH^{-1}
  =I+aE_{12}(I+C)^{-1}
  =I+\sum_{k=2}^{r}(-1)^{k-2}
      a\left(\prod_{j=2}^{k-1}c_j\right)E_{1k}.
\end{equation}
The empty product for $k=2$ is $1$. Every nonconstant coefficient of $P$ is a negative Laurent part, while $H$ is holomorphic. Hence \eqref{eq:chainP} is the normalized polar factorization, by uniqueness. All signs follow from the finite geometric inverse of $I+C$.

\subsection{The exponent design}

For $r\geq4$ and $n\geq2$, set
\begin{equation}\label{eq:deepdata}
 a_n=\frac{t}{z-n},\qquad
 c_{2,n}=t^{1-1/n},\qquad
 c_{j,n}=t\ (3\leq j\leq r-2),\qquad
 c_{r-1,n}=t^{2/n}.
\end{equation}
The middle range is empty for $r=4$. Let $C_n$ and $G_n$ be the matrices just defined from these data. Every $G_n$ is finite, upper unitriangular, and has just one raw singular matrix coefficient, namely $E_{12}/(z-n)$ at exponent $1$.

For $3\leq k\leq r-1$, the exponent in $(P_n)_{1k}$ is
\begin{equation}\label{eq:lowerexponents}
 \lambda_{k,n}=k-1-\frac1n,
\end{equation}
whereas the exponent in the final entry is
\begin{equation}\label{eq:lastexponent}
 \lambda_{r,n}=r-2+\frac1n.
\end{equation}
For each fixed $k<r$, the exponents \eqref{eq:lowerexponents} increase with $n$, and their set has order type $\omega$. There are finitely many such sets, so their union with $\{1\}$ is well ordered. In contrast, the final exponents \eqref{eq:lastexponent} strictly decrease.

\begin{theorem}[Arbitrarily deep central support obstructions]\label{thm:depth}
For each integer $r\geq3$ there exists a positive upper-unitriangular cocycle on the puncture-and-disk cover of $\C$ with $A=\{2,3,\ldots\}$, over $\Gamma=\Q$, having all the following properties:
\begin{enumerate}[label=\textup{(\roman*)}]
\item Its original singular coefficient support is $\{1\}$ and its determinant is one.
\item Its image in the quotient of the upper-unitriangular group by $Z_r=\{I+cE_{1r}\}$ is a coboundary.
\item It is not a coboundary in the full group of positive $r\times r$ matrices.
\item It remains nonsplit after every ordered value-group extension.
\end{enumerate}
For $r=3$ use \eqref{eq:badG}; for $r\geq4$ use \eqref{eq:deepdata} and the chain construction.
\end{theorem}

\begin{proof}
The first property was checked directly. Formula \eqref{eq:chainP} and \eqref{eq:lastexponent} show that the normalized polar support contains a descending infinite sequence. Thus \cref{thm:criterion} proves (iii), and \cref{cor:extension} proves (iv).

For (ii), let $P_n'$ be obtained from $P_n$ by deleting its $(1,r)$ entry. Their global support union is well ordered by \eqref{eq:lowerexponents}; the case $r=3$ leaves only exponent $1$. Every $P_n'-I$ belongs to the row-one algebra
\[
 V=\operatorname{span}\{E_{12},\ldots,E_{1,r-1}\},
 \qquad V^2=0.
\]
The gluing problem for the $P_n'$ is therefore additive and splits by the scalar construction with this one well-ordered support union. Explicitly, solve the ordinary principal-part problem for each coefficient of each allowed row entry. This gives $F'$ on $U_0$ and $Q_n'$ on $D_n$, valued in $I+V$, such that $F'=P_n'Q_n'$.

In the quotient by $Z_r$, one has $\overline P_n=\overline P_n'$. Since $G_n=P_nH_n$ with $H_n$ holomorphic, the identities
\[
 \overline G_n
 =\overline F'\,
   \bigl(\overline H_n^{-1}\overline Q_n'\bigr)^{-1}
\]
provide the required coboundary. All factors have individually legitimate positive supports. No matrix representative for the quotient is silently assumed; the bars denote the ordinary group quotient by the central subgroup.
\end{proof}

Here $Z_r$ means the central subgroup with $c\in\JJ$. Over $\Gamma=\Q$, one has $\JJ^m=\JJ$ for every positive integer $m$: a nonzero section of least exponent $\delta>0$ factors into $m-1$ monomials $t^{\delta/m}$ and one remaining positive section. The elementary commutator identity $[I+aE_{ij},I+bE_{jk}]=I+abE_{ik}$ then shows that $Z_r$ is the final nonzero term of the lower central series of the positive upper-unitriangular group. Every further lower-central quotient factors through the quotient in (ii). Thus all those lower-depth tests pass. The full obstruction can first appear at depth $r-1$, for arbitrarily large finite $r$.

\begin{remark}[What the depth assertion means]
The theorem does not say that a scalar equation is incapable of detecting the final obstruction once a noncommutative normal form has been found. The final coefficient is, of course, scalar. It says that discarding the last central level before normalizing loses an obstruction created by the multiplication law. The new information is located in a genuinely higher noncommutative compatibility condition.
\end{remark}

\section{Integral vector bundles and local invisibility}\label{sec:bundles}

\subsection{From positive matrices to actual vector bundles}

The transition matrices $G_a$ define a locally free rank-$r$ sheaf $\mathcal E_G$ over $\II$ on the ordinary surface $X$. Its ordinary reduction is canonically the trivial rank-$r$ holomorphic bundle, because every transition reduces to $I$. A positive splitting is a trivialization compatible with that reduced framing.

\begin{proposition}[For triviality, the reduced framing costs nothing]\label{prop:integraltrivial}
For a bundle defined by positive transition matrices, the following are equivalent:
\begin{enumerate}[label=\textup{(\roman*)}]
\item $\mathcal E_G$ is trivial as an $\II$-module.
\item Its transition cocycle splits by positive matrices.
\end{enumerate}
The implication holds on an arbitrary ordinary complex base, not only on a star cover.
\end{proposition}

\begin{proof}
A positive splitting gives a trivialization. Conversely, suppose matrices $A_i\in\GL_r(\II(U_i))$ satisfy $G_{ij}=A_iA_j^{-1}$. Reduce the equality modulo $\JJ$. Since $\red G_{ij}=I$, the ordinary invertible matrices $\red A_i$ agree on overlaps and glue to $A\in\GL_r(\OO(X))$. Multiplication on the right by the ordinary matrix $A^{-1}$ produces $B_i=A_iA^{-1}$ with reduction $I$. Hence $B_i\in\mathscr G_{r,\Gamma}(U_i)$ and $G_{ij}=B_iB_j^{-1}$.
\end{proof}

\begin{corollary}[Nontrivial determinant-one integral bundles]\label{cor:nontrivialbundles}
The examples in \cref{thm:hidden,thm:depth} define nontrivial locally free $\II$-modules even after forgetting the upper-triangular flag and the chosen reduced framing. Their ordinary reductions and determinant line bundles are trivial.
\end{corollary}

\begin{proof}
A general integral trivialization would give a positive matrix splitting by \cref{prop:integraltrivial}, contradicting \cref{thm:criterion}. Their transition determinants are exactly one, not merely one modulo positive exponents.
\end{proof}

\subsection{Every bounded restriction is trivial}

\begin{corollary}[Exhaustion-invisible nontriviality]\label{cor:bounded}
Each of the integral bundles in \cref{cor:nontrivialbundles} restricts trivially to every bounded ordinary domain in $\C$. In particular it is trivial on every disk of an increasing exhaustion of the plane, but not on the plane itself.
\end{corollary}

\begin{proof}
A bounded domain meets only finitely many of the fixed disks $D_n$. On the restricted cover there are therefore only finitely many nonidentity transition matrices, each with legitimate positive support. Their union is well ordered. The support-controlled Stein splitting theorem proved in \cref{thm:Stein} below applies, since every ordinary plane domain is Stein. Equivalently, apply the finite-puncture construction directly and absorb transitions on pieces that contain no puncture into holomorphic local gauges.
\end{proof}

The failure is compatibility of the infinitely many local trivializations, not existence on any one bounded domain. Consequently a numerical experiment involving finitely many punctures cannot establish global triviality. It will necessarily miss these examples.

\subsection{The localization boundary}

Let $\HH=\OO((t^\Gamma))$ denote the sheaf in which arbitrary monomial shifts are allowed. The natural map $\II\to\HH$ extends scalars from $\mathcal E_G$ to a rank-$r$ $\HH$-module. We have \emph{not} proved that the resulting module is nontrivial. The reduction argument in \cref{prop:integraltrivial} cannot be repeated for matrices over $\HH$: there is no reduction ring homomorphism $\HH\to\OO$ killing positive monomials, because those monomials are units in $\HH$.

Accordingly, the precise conclusions are nonexistence of a positive coherent matrix splitting and nontriviality of the corresponding \emph{integral} coherent vector bundle. Whether negative-valuation changes of frame can trivialize particular higher-rank examples is a separate question. The scalar report proves more in rank one by a special leading-unit normalization \cite{RepoGlobal}; no matrix analogue of that argument is assumed here.

This distinction matters for the interpretation as surcomplex analysis. Integral models describe holomorphic matrices with a defined ordinary reduction and infinitesimal corrections. They are not the entire unrestricted matrix geometry over the field $\No[i]$. Our theorem is strong within the former category and intentionally does not conflate the two.

\section{Positive rigidity on Stein manifolds}\label{sec:Stein}

\subsection{An arbitrary-cover theorem with a support hypothesis}

Let $X$ be a connected Stein manifold and $(U_i)_{i\in I}$ a set-indexed open cover. On disconnected intersections, support statements refer to all connected components. Suppose
\[
 G_{ij}\in\mathscr G_{r,\Gamma}(U_i\cap U_j),
 \qquad G_{ij}G_{jk}=G_{ik},\quad G_{ii}=I.
\]
Call the cocycle \emph{globally support-admissible} if
\begin{equation}\label{eq:Ecover}
 E=\bigcup_{i,j}\supp(G_{ij}-I)
\end{equation}
is well ordered. This is a sufficient condition on a presentation, not an invariant necessary condition; regular local gauges may destroy it without destroying triviality.

\begin{theorem}[Support-controlled Stein splitting]\label{thm:Stein}
A globally support-admissible positive matrix cocycle on a Stein manifold splits on its original cover:
\[
 G_{ij}=B_iB_j^{-1},
 \qquad B_i\in\mathscr G_{r,\Gamma}(U_i).
\]
All supports $\supp(B_i-I)$ can be contained in the one monoid $\angles E\setminus\{0\}$.
\end{theorem}

\begin{proof}
Set $M=\angles E$. We solve $G_{ij}B_j=B_i$ coefficientwise, with
\[
 B_i=I+\sum_{\gamma\in M\setminus\{0\}}b_{i,\gamma}t^\gamma.
\]
Assume the equations hold at all exponents below $\gamma$. The coefficient equation at $\gamma$ is
\begin{equation}\label{eq:Steinrec}
 b_{i,\gamma}-b_{j,\gamma}
 =g_{ij,\gamma}+
  \sum_{\substack{\eta+\theta=\gamma\\\eta,\theta>0}}
     g_{ij,\eta}b_{j,\theta}
 =:R_{ij,\gamma}.
\end{equation}
The sum is finite and uses only earlier coefficients. The identity $G_{ij}G_{jk}=G_{ik}$, together with the equations already solved below $\gamma$, implies
\[
 R_{ij,\gamma}+R_{jk,\gamma}=R_{ik,\gamma}
\]
on triple overlaps. One may check this by expanding $G_{ij}(G_{jk}B_k-B_j)+(G_{ij}B_j-B_i)=G_{ik}B_k-B_i$ at exponent $\gamma$; all positive multipliers of a lower defect vanish by induction.

Thus $R_{ij,\gamma}$ is an ordinary additive cocycle with values in $\OO^{r^2}$. Cartan's theorem B gives $H^1(X,\OO^{r^2})=0$, so it has a coboundary $b_{i,\gamma}$ on the original cover. This is again the torsor interpretation of first cohomology, and does not assume that the Hahn sheaf is coherent. Recursion on $M\setminus\{0\}$ gives the required matrices with a common support bound. Positive inversion proves that they are invertible. No new ordinary domains and no new Hahn exponents are introduced.
\end{proof}

\subsection{Deformations of a nontrivial ordinary bundle}

The same principle is not limited to a trivial ordinary reduction. Let $\mathcal V$ be an ordinary holomorphic rank-$r$ bundle on $X$, with transition matrices $g^0_{ij}$. Suppose Hahn transition matrices $G_{ij}$ have reduction $g^0_{ij}$ and are invertible over $\II$. The relative transitions $G_{ij}(g^0_{ij})^{-1}-I$ have positive support on overlaps.

\begin{theorem}[Support-controlled deformation rigidity]\label{thm:deformation}
If the union of these relative positive supports is well ordered and $X$ is Stein, there exist positive local matrices $B_i$ such that
\begin{equation}\label{eq:deformation}
 G_{ij}=B_i g^0_{ij} B_j^{-1}.
\end{equation}
The nonconstant supports of all $B_i$ lie in the monoid generated by the given support union. Thus the integral Hahn deformation is isomorphic, by an isomorphism reducing to the identity, to the scalar extension of its ordinary bundle.
\end{theorem}

\begin{proof}
Proceed as in \cref{thm:Stein}, now solving $G_{ij}B_j=B_i g^0_{ij}$. At a fixed positive exponent, multiply the residual equation on the right by $(g^0_{ij})^{-1}$. The unknown coefficients appear as
\[
 b_{i,\gamma}-g^0_{ij}b_{j,\gamma}(g^0_{ij})^{-1}.
\]
The residual is an additive cocycle for the ordinary endomorphism sheaf $\End(\mathcal V)$; the cocycle identity follows from $G_{ij}G_{jk}=G_{ik}$ and the lower solved equations exactly as before. This sheaf is locally free, hence coherent. Cartan's theorem B gives $H^1(X,\End(\mathcal V))=0$. Solve the ordinary coefficient cocycle and continue the support-controlled recursion. All multiplications by $g^0_{ij}$ and its inverse occur at exponent zero and do not enlarge the exponent set.
\end{proof}

This is a formal deformation-rigidity argument with a precise support certificate. It should not be interpreted as a new proof of the ordinary Oka--Grauert principle, nor as an unconditional theorem about every Hahn-valued deformation on a noncompact base.

\subsection{A cyclic versus noncyclic dichotomy}

\begin{lemma}[Ordered-group alternative]\label{lem:groups}
For a nonzero ordered abelian group $\Gamma$, the positive cone contains no infinite strictly decreasing sequence if and only if $\Gamma=\Z\delta$ for a least positive element $\delta$.
\end{lemma}

\begin{proof}
The cyclic case is immediate. Conversely, if there is no least positive element, choose a smaller positive element repeatedly. If there is a least positive element $\delta$ but the group is not $\Z\delta$, choose a positive $\gamma\notin\Z\delta$. Then $\gamma-n\delta>0$ for every integer $n\geq0$: otherwise the first nonpositive difference would leave an element strictly between $0$ and $\delta$, or imply $\gamma\in\Z\delta$. The sequence $\gamma,\gamma-\delta,\ldots$ is strictly decreasing and positive.
\end{proof}

\begin{theorem}[Universal positive rigidity and its failure]\label{thm:dichotomy}
Fix a rank $r\geq1$. Every positive rank-$r$ cocycle on every Stein manifold is a coboundary if and only if
\[
 \Gamma=\{0\}\quad\text{or}\quad\Gamma=\Z\delta\ \text{for some }\delta>0.
\]
For every noncyclic nonzero $\Gamma$ and every $r\geq2$, there is a nonsplit determinant-one positive rank-$r$ cocycle already on the puncture-and-disk cover of $\C$.
\end{theorem}

\begin{proof}
In the cyclic case every transition support is a subset of the universal well-ordered positive set $\{\delta,2\delta,\ldots\}$. Its global union is automatically well ordered, so \cref{thm:Stein} applies on any cover. The zero group gives only the identity cocycle.

For the converse, choose a strictly decreasing positive sequence $\gamma_n$ by \cref{lem:groups}, and set
\[
 g_n=\exp\left(\frac{t^{\gamma_n}}{z-n}\right).
\]
This scalar matrix is already its own normalized polar factor: every nonconstant coefficient is a multiple of a negative power of $z-n$. Its support union contains the descending set $\{\gamma_n\}$, so it is nonsplit. In rank $r\geq2$, use
\begin{equation}\label{eq:SLcounter}
 G_n=\operatorname{diag}(g_n,g_n^{-1},1,\ldots,1).
\end{equation}
Every nonconstant coefficient of this matrix is polar, so $P_n=G_n$. The same descending support prevents splitting in the full positive matrix group, while $\det G_n=1$. For rank one, the scalar example suffices. \Cref{prop:integraltrivial} also converts these examples to nontrivial integral vector bundles.
\end{proof}

The rank-one part agrees with the scalar dichotomy in \cite{RepoGlobal}. The higher-rank conclusion is stated for integral positive deformations, and the rank-three and arbitrary-depth examples provide the nonabelian information that a diagonal example by itself would not supply.

\section{How to compute the invariant, and what computation cannot decide}\label{sec:computation}

\subsection{A support-aware procedure}

For a presented star-cover datum, the mathematical procedure is:
\begin{enumerate}
\item At each puncture, compute the normalized factor by \eqref{eq:factorrecursion1}--\eqref{eq:factorrecursion2}, or exploit a finite matrix identity such as \eqref{eq:HeisenbergP} or \eqref{eq:chainP}.
\item Extract the union of the supports of the normalized polar coefficients, not the supports of the raw entries and not just the leading valuations.
\item Supply a proof of well ordering or a descending-sequence obstruction for that union. In the admissible case use \eqref{eq:globalrecursion}--\eqref{eq:qrecursion} to construct a splitting with a monoid support certificate.
\end{enumerate}
These steps give an exact criterion, not a decision algorithm for an arbitrary black-box infinite sequence of functions. In general, an input description must support a proof of the required infinite order property. Merely checking many exponents or many punctures does not settle it.

\subsection{Finite truncations}

For finitely generated positive rational exponent monoids, every bounded exponent interval is finite. A sparse dictionary keyed by exact rational exponents therefore suffices for a finite-cutoff implementation of the local recursion. Each dictionary value can be a matrix of ordinary Laurent polynomials. The projection $\pi_-$ then keeps the negative powers of $u$ and $\pi_+$ keeps the nonnegative powers.

For a general higher-rank ordered group, a bounded interval can be infinite, and no such finite-cutoff assumption is justified. A practical implementation must either work with a certified finite dependency set or restrict its representation class. This is the computational counterpart of \cref{rem:nonarch}.

\subsection{The accompanying exact checks}

The archive includes \texttt{code/verify.py}, using Python and SymPy with exact rational arithmetic. It checks the rank-three factors and the missing cross term, the repaired example, the chain factorization for several ranks, the noncommuting local coefficient recursion at a rational-exponent cutoff, invariance under a holomorphic right gauge at that cutoff, and finite-puncture splittings. It also records finite instances of the exponent formulas \eqref{eq:lowerexponents}--\eqref{eq:lastexponent}.

The recorded run used Python 3.13.5 and SymPy 1.14.0; all 2,899 exact checks passed. The output is in \texttt{data/verification.txt}. These checks audit algebraic signs, multiplication order, and finite coefficient identities. They do not prove the support lemma, Cartan's theorem B, an infinite descending obstruction, or novelty. The general proofs are the arguments in the body of this article.

\section{Relation to the literature and remaining questions}\label{sec:literature}

\subsection{Three distinct ingredients}

The first ingredient is generalized-series algebra. Neumann's positive-support lemma supplies well ordering and finite multiplicity \cite{Neumann}. For surreal fields, the analytic interpretation of restricted ordinary functions in suitable Hahn workspaces is established in \cite{vdDE}; the broader landscape is surveyed in \cite{MantovaMatusinski}. Our use is modest: complex coefficients, set-sized supports, ordinary holomorphic coefficient functions, and positive Taylor evaluation. No scalar derivation on $\No$ is chosen.

The second ingredient is normalized matrix factorization. The noncommutative Birkhoff construction in \cite{EFGK}, especially its recursive factorization of $1+$ filtered algebras with an idempotent splitting operator, is a close algebraic antecedent. The Laurent decomposition gives exactly the relevant pair of subalgebras. The local theorem here supplies a direct Hahn-word argument for arbitrary ordered $\Gamma$ and keeps every ordinary coefficient on a fixed disk. We do not claim that replacing a complete filtered algebra by an explicit summable Hahn presentation creates a new general factorization idea. Classical analytic Birkhoff and bundle-splitting theory is discussed in \cite{Wone}; its usual single-loop setting is different from the infinitely many independently supported gluing data treated here.

The third ingredient is ordinary analytic gluing. On the plane it is supplied by the explicit construction in \cref{lem:MLplane}; on an open surface or a Stein manifold it is supplied by Behnke--Stein and Cartan B \cite{BS,Cartan,Schmitt}. The crucial point is that these results are applied at \emph{one complex coefficient at a time}. They do not by themselves ensure that the outputs across exponents assemble into a valid Hahn section. The exact polar-support criterion isolates that missing condition.

\subsection{The novelty assessment}

The targeted check examined the repository's global-support report and its explicitly listed future directions, generalized-series sources, and primary papers on algebraic noncommutative Birkhoff factorization. It located established local factorization machinery but did not locate the equivalence \eqref{eq:criterionintro} for independently supported matrix data on a countable puncture-and-disk cover, or the arbitrary-depth examples of \cref{thm:depth}. Some broad web searches returned irrelevant material; those results were not treated as evidence of absence. The novelty claim is therefore limited to these precise formulations and examples, to the best of the present literature check.

The article resolves a specific finite-rank support problem raised by the repository, under the stated positive integral and star-cover hypotheses. It is not a declaration that the full higher-rank theory or an established named conjecture has been settled.

\subsection{Questions exposed by the results}

\paragraph{Monomial localization.}
Does the extension $\mathcal E_G\otimes_{\II}\HH$ remain nontrivial for the explicit hidden and deep examples? A proof would need control of general Hahn matrix gauges, not just positive ones. The ordinary reduction normalization is unavailable after localization. This question is intentionally not answered here.

\paragraph{Arbitrary-cover normal forms.}
The polar support is exact on the puncture-and-disk cover because every overlap has a distinguished isolated Laurent singularity. An equally explicit invariant for general covers would require a substitute for this canonical local negative factor. The sufficient globally admissible-support theorem does not provide that substitute.

\paragraph{Classification beyond triviality.}
The criterion decides whether a given positive star-cover cocycle represents the trivial integral bundle. A classification of all isomorphism classes would require understanding the action of global and local gauges on pairs of polar families, not just testing one family against the identity. The right-torsor statement describes all solutions once a splitting exists, not all nonsplit objects.

\paragraph{Higher-dimensional singular sets.}
For a Stein manifold of dimension greater than one, the support-controlled rigidity theorem still applies. The sharp star-cover invariant does not immediately generalize: there is no one-variable Laurent projection for an arbitrary positive-dimensional analytic singular locus, and its ideal-theoretic and cohomological replacement may have nontrivial relations.

\paragraph{Formal verification.}
The local support recursion, chain identity, and necessity direction appear well suited to formalization using Hahn series and finite matrix algebra. The full global sufficiency proof additionally needs a formal ordinary Mittag--Leffler theorem or a suitable sheaf-cohomology interface. The accompanying script is not a substitute for such a formalization.

\section{Conclusion}

The correct obstruction to positive matrix gluing is not the support of the input entries, nor the support of their raw principal parts, nor a determinant. One must first remove the holomorphic factors in the correct noncommutative order. The resulting normalized polar support gives an exact necessary and sufficient criterion on a puncture-and-disk cover.

The explicit examples show why this normalization is mathematically consequential. A family with one singular input exponent can force a descending support at a hidden matrix entry. Adding a cross term can repair the family while worsening its raw singular support. The same phenomenon can be postponed to the final central layer of upper-unitriangular groups of arbitrarily large finite rank. These yield nontrivial determinant-one integral Hahn bundles with trivial ordinary reduction and trivial bounded restrictions.

At the opposite boundary, a single globally well-ordered support certificate restores Stein rigidity in every finite rank. Cyclic value groups supply such a certificate automatically; noncyclic groups do not. This separates ordinary analytic solvability, noncommutative compatibility, and generalized-series admissibility into three precise and independently visible parts.

\appendix

\section{A proof and support audit}\label{app:audit}

\begin{longtable}{@{}L{0.25\textwidth}L{0.35\textwidth}L{0.32\textwidth}@{}}
\toprule
Construction & Support control & Ordinary analytic requirement\\
\midrule
\endhead
Positive matrix inverse & Positive monoid generated by the input support; finitely many exponent words and finite matrix paths at each coefficient & Coefficients remain on the original domain.\\[0.4em]
Local factorization & $\angles{\supp(G-I)}$; transfinite triangular recursion & Laurent splitting on one fixed punctured disk, with no shrinking.\\[0.4em]
Necessity of the global criterion & All $P_a-I$ lie in the one monoid generated by $\supp(B_0-I)$ & Right-gauge invariance and uniqueness of local factorization.\\[0.4em]
Global construction & $\angles{\EU(G)}$ for the central matrix and the $Q_a$ & One ordinary isolated-singularity interpolation problem per matrix coefficient and exponent.\\[0.4em]
Disk trivializers & Generated monoid of the global polar support and that disk's original support & Products of matrices holomorphic on that disk.\\[0.4em]
Right torsor & Finite union of the two global solution supports and positive monoid closure & Identity-theorem extension across each puncture.\\[0.4em]
Stein splitting & One monoid generated by all transition supports & Ordinary $H^1(X,\OO^{r^2})=0$.\\[0.4em]
Deformation rigidity & The same positive monoid; ordinary transition matrices are exponent-zero coefficients & Ordinary $H^1(X,\End(\mathcal V))=0$.\\[0.4em]
Rank-three obstruction & Forced descending set $\{1+1/n\}$ & Only simple poles are used.\\[0.4em]
Depth-$r-1$ obstruction & Lower partial-sum supports are well ordered; final support $\{r-2+1/n\}$ descends & Only simple poles and holomorphic constants are used.\\[0.4em]
Integral nontriviality & A hypothetical integral frame normalizes to a positive frame & Reduction of invertible matrices and gluing of the reduced frame.\\
\bottomrule
\end{longtable}

The indispensable hypotheses are as follows. The value group is a set and is totally ordered. Supports are well ordered, not merely bounded below. Matrix rank is finite. The infinite puncture set is ordinary closed discrete. Coefficient functions share their stated ordinary domains. The gluing matrices reduce to the identity, or to a specified ordinary bundle in \cref{thm:deformation}. None of these hypotheses is replaced by fine-topological convergence.

\section{A dependency map}

The main proof chain is
\[
\begin{gathered}
\text{Neumann support lemma}+\text{ordinary Laurent splitting}\\
\Downarrow\\
\text{local normalized factorization and right-gauge invariance}\\
\Downarrow\quad\text{with ordinary isolated-singularity interpolation}\\
\text{exact global polar-support criterion}\\
\Downarrow\\
\text{rank-three cross-term obstruction and arbitrary-depth examples}\\
\Downarrow\quad\text{with reduction of an integral frame}\\
\text{nontrivial integral vector bundles.}
\end{gathered}
\]
The arbitrary-cover Stein theorem is a separate positive branch. Its only cohomological input is ordinary coherent-sheaf acyclicity. It is not used to justify the infinite star-cover criterion on the plane, whose interpolation step has its own direct proof. It is used for the clean statement about all bounded restrictions and for the universal cyclic-group result.

\section{Reproduction and repository provenance}

The supplied \texttt{article.tex} has an internal bibliography and no external graphics or font files. A standard TeX Live installation can build it with
\begin{verbatim}
latexmk -pdf -interaction=nonstopmode -halt-on-error article.tex
\end{verbatim}
The verification script runs with
\begin{verbatim}
python code/verify.py
\end{verbatim}
A \texttt{README.md} records the exact scope and the code dependencies. The repository was read at the pinned commit stated in \cref{sec:introduction}; no changes were pushed to it. This article is a new standalone continuation, not a replacement for the repository's existing merged reports.

\begin{thebibliography}{99}
\small\raggedright
\setlength{\itemsep}{0.3em}

\bibitem{RepoMap}
V.~Reshetnikov, \emph{Surreal repository: report catalogue},
\texttt{docs/README.md}, commit
\texttt{39f2be6667ade51bca2b45daa47e289d69c09764}, inspected September 21, 2026.
\url{https://github.com/VladimirReshetnikov/Surreal/blob/39f2be6667ade51bca2b45daa47e289d69c09764/docs/README.md}.

\bibitem{RepoGlobal}
\emph{Global Support Obstructions: Moving Divisors, Mittag--Leffler, and a Picard Dichotomy}, AI-assisted merged research draft in V.~Reshetnikov's \emph{Surreal} repository, September 21, 2026. Pinned commit as above; especially ``Further precise directions.''
\url{https://github.com/VladimirReshetnikov/Surreal/blob/39f2be6667ade51bca2b45daa47e289d69c09764/docs/surcomplex/global-divisors/article.tex}.

\bibitem{Neumann}
B.~H.~Neumann, On ordered division rings,
\emph{Transactions of the American Mathematical Society} \textbf{66} (1949), 202--252.
\url{https://www.jstor.org/stable/1990552}.

\bibitem{vdDE}
L.~van den Dries and P.~Ehrlich, Fields of surreal numbers and exponentiation,
\emph{Fundamenta Mathematicae} \textbf{167} (2001), 173--188.
\url{https://doi.org/10.4064/fm167-2-3}.

\bibitem{MantovaMatusinski}
V.~Mantova and M.~Matusinski,
\emph{Surreal numbers with derivation, Hardy fields and transseries: a survey}, arXiv:1608.03413.
\url{https://arxiv.org/abs/1608.03413}.

\bibitem{EFGK}
K.~Ebrahimi-Fard, L.~Guo, and D.~Kreimer,
Spitzer's identity and the algebraic Birkhoff decomposition in pQFT,
\emph{Journal of Physics A: Mathematical and General} \textbf{37} (2004), 11037--11052. See especially Theorems 3.6--3.7 of the preprint.
\url{https://arxiv.org/abs/hep-th/0407082}.

\bibitem{EFMP}
K.~Ebrahimi-Fard, D.~Manchon, and F.~Patras,
\emph{A noncommutative Bohnenblust--Spitzer identity for Rota--Baxter algebras solves Bogoliubov's recursion}, arXiv:0705.1265.
\url{https://arxiv.org/abs/0705.1265}.

\bibitem{BS}
H.~Behnke and K.~Stein, Entwicklung analytischer Funktionen auf Riemannschen Fl\"achen,
\emph{Mathematische Annalen} \textbf{120} (1947/1949), 430--461.
\url{https://doi.org/10.1007/BF01447838}.

\bibitem{Cartan}
H.~Cartan, Vari\'et\'es analytiques complexes et cohomologie,
in \emph{Colloque sur les fonctions de plusieurs variables}, Brussels, 1953, 41--55. Classical source for the coherent-sheaf acyclicity theorem on Stein manifolds.

\bibitem{Schmitt}
M.~M.~Schmitt,
\emph{Kohomologie mit Schranken und Fortsetzung holomorpher Funktionen durch lineare stetige Operatoren}, doctoral thesis, arXiv:math/0202263 (2002).
Includes a proof of Cartan's theorem B for coherent subsheaves of finite free holomorphic sheaves and coefficientwise coboundary constructions.
\url{https://arxiv.org/abs/math/0202263}.

\bibitem{Wone}
O.~Wone, \emph{The Birkhoff--Grothendieck theorem}, arXiv:2304.08037, version 2.
Historical and mathematical exposition of bundle splitting and Birkhoff factorization.
\url{https://arxiv.org/abs/2304.08037}.

\end{thebibliography}
\end{document}
